\documentclass[11pt,a4paper]{article}

\usepackage[margin=1.05in]{geometry}
\usepackage{amsmath,amssymb,amsthm}
\usepackage{mathtools}
\usepackage{stmaryrd}
\usepackage{microtype}
\setlength{\emergencystretch}{2.5em}
\usepackage{enumitem}
\usepackage{booktabs}
\usepackage{array}
\usepackage{xcolor}
\usepackage[colorlinks=true,linkcolor=blue!45!black,urlcolor=blue!60!black]{hyperref}

\newtheorem{theorem}{Theorem}[section]
\newtheorem{proposition}[theorem]{Proposition}
\newtheorem{lemma}[theorem]{Lemma}
\newtheorem{corollary}[theorem]{Corollary}
\newtheorem{conjecture}[theorem]{Conjecture}
\theoremstyle{definition}
\newtheorem{definition}[theorem]{Definition}
\newtheorem{example}[theorem]{Example}
\theoremstyle{remark}
\newtheorem{remark}[theorem]{Remark}
\newtheorem{refuted}[theorem]{Refuted conjecture}
\newtheorem{question}[theorem]{Question}

\newcommand{\FF}{\mathbb{F}}
\newcommand{\NN}{\mathbb{N}}
\newcommand{\ZZ}{\mathbb{Z}}
\newcommand{\QQ}{\mathbb{Q}}
\newcommand{\cat}[1]{\mathbf{#1}}
\newcommand{\Set}{\cat{Set}}
\newcommand{\Vect}{\cat{Vect}}
\newcommand{\Rep}{\mathrm{Rep}}
\newcommand{\ext}[1]{\llbracket #1 \rrbracket}
\newcommand{\con}{\triangleleft}
\DeclareMathOperator{\Hom}{Hom}
\DeclareMathOperator{\GL}{GL}
\DeclareMathOperator{\End}{End}
\DeclareMathOperator{\Ind}{Ind}
\newcommand{\qbin}[2]{\genfrac{[}{]}{0pt}{}{#1}{#2}_{q}}
\newcommand{\Fl}{\mathrm{Fl}}
\newcommand{\Schub}{\mathfrak{S}}

\definecolor{okgreen}{rgb}{0.0,0.45,0.15}
\definecolor{nored}{rgb}{0.65,0.0,0.0}
\newcommand{\OK}{\textcolor{okgreen}{\textbf{confirmed}}}
\newcommand{\NO}{\textcolor{nored}{\textbf{refuted}}}

\title{\bfseries The Flag Container\\[0.4em]
\large why it gives Demazure operators and not Schubert polynomials}
\author{Robin Langer \and Claude}
\date{20 September 2026}

\begin{document}
\maketitle

\begin{abstract}
The Targets table of \texttt{q-container.pdf} has rows for the $q$-Leibniz rule,
$q$-Pascal, Pieri--Littlewood--Richardson and Macdonald polynomials, and no row
for Schubert polynomials. This note attempts to earn one and instead establishes
that the row cannot be written in the form expected.

Schubert polynomials are generated from the top class by divided differences
$\partial_i f = (f - s_i f)/(x_i - x_{i+1})$, and the container derivative of
\texttt{q-container.pdf} Theorem~D is the Jackson operator
$D_q f = (f - \sigma_q f)/((1-q)x)$. Both have the shape
(\emph{$1$ minus a group element}) over (\emph{a difference of weights}), Jackson
dividing by a scaling and Demazure by a transposition; and a container already
carries transpositions, since relabelling positions is what a container's symmetry
\emph{is}. The conjecture was therefore that $\partial_i$ needs no new machinery.

Half of that survives. Building the operator from the fibration that forgets the
$i$-th subspace of a flag, we verify by exact enumeration
($n \le 4$, $q = 2, 3$) that it descends to the cell module and satisfies the
\emph{braid} relations, so $\Schub_w$ is independent of the reduced word at
container level with no new machinery, as predicted. But the operator satisfies
the \emph{Hecke} quadratic and not the nilCoxeter relation $\partial_i^2 = 0$;
and we prove (Theorem~\ref{thm:nonil}) that no operator attached to a single
index can be nilpotent unless $q = -1$. The container hands you an idempotent
where Schubert polynomials need a nilpotent, and an idempotent divided difference
is a \emph{Demazure} operator, whose polynomials are key polynomials.

The consequence is a correction to the dictionary rather than an addition to it:
counting $\FF_q$-points sees the Hecke algebra and the Poincar\'e polynomial
$\sum_w q^{\ell(w)} = [n]_q!$, while divided differences live in the Chow ring,
which is a different decategorification of the same cells.

\medskip
\noindent\emph{Companion:} \texttt{q-container.pdf}; its remark that the Hecke algebra is the endomorphism algebra of
the flag container is the starting point. Code: \texttt{code/flag.py},
results recorded in \texttt{code/FINDINGS.md}.
\end{abstract}

\section{The question, and why it is not idle}\label{sec:q}

\texttt{q-container.pdf} observes that $\Ind_B^G \mathbf{1} = k[G/B]$ is a free
module on a hom-set --- a container in the $k[-]$ sense --- and that
$\End_G(\Ind_B^G \mathbf{1}) = H_n(q)$, so that \emph{the Hecke algebra is the
endomorphism algebra of the flag container}. Since the Schubert basis of
$H^*(\Fl_n)$ is indexed by the same Weyl group that indexes $H_n(q)$, it is
natural to expect the Schubert calculus to be visible from the container side.

The expectation has a precise form. Write $\partial_i$ for the divided difference
and $D_q$ for the Jackson $q$-derivative of Theorem~D:
\[
  \partial_i f = \frac{f - s_i f}{x_i - x_{i+1}},
  \qquad\qquad
  D_q f = \frac{f - \sigma_q f}{(1-q)\,x}, \quad \sigma_q f(x) = f(qx).
\]
Both are \emph{divided differences for a group element}: a difference operator
over the corresponding difference of weights. Theorem~D says the container
derivative is the Jackson one, where the group element is the scaling $\sigma_q$.
Containers carry the symmetric group on positions for free, so nothing new should
be needed to build the operator with $s_i$ in place of $\sigma_q$.

\begin{conjecture}[under test]\label{conj:nil}
The flag container carries operators $\partial_i$, built from its own structure,
generating the nilCoxeter algebra: $\partial_i^2 = 0$, together with the braid
relations, so that $\Schub_w = \partial_{w^{-1}w_0}(x^\delta)$ is well defined.
\end{conjecture}

\begin{remark}[what would make this circular]
Defining $\partial_i$ by its action on Schubert classes --- $\partial_i \Schub_w
= \Schub_{ws_i}$ when $\ell(ws_i) < \ell(w)$ and $0$ otherwise --- and then
verifying the nilCoxeter relations would verify nothing. Below the operator is
built from the geometry and its relations are read off.
\end{remark}

\section{The flag container}

\begin{definition}[flag container]\label{def:flag}
Let $\Fl_n(\FF_q)$ be the set of complete flags $0 = F_0 < F_1 < \cdots < F_n =
\FF_q^n$, and $E$ the standard flag. The \emph{relative position} of $F$ to $E$
is the permutation $w \in S_n$ determined by the rank matrix
$r_{ij} = \dim(E_i \cap F_j)$ via
\[
  w(j) = i \iff r_{ij} - r_{i-1,j} - r_{i,j-1} + r_{i-1,j-1} = 1 .
\]
The \emph{flag container} is $S_n \con \bigl(w \mapsto \FF_q^{\ell(w)}\bigr)$,
with extension $\ext{\,\cdot\,}X = \sum_{w \in S_n} X^{\ell(w)}$.
\end{definition}

So the shapes are the permutations and the positions over shape $w$ are the cell
$C_w$. This is a $q$-container in the sense of \texttt{q-container.pdf}
Definition~2: the shapes form a set, the positions are $\FF_q$-spaces.

\begin{proposition}[Bruhat, by enumeration]\label{prop:bruhat}
$|\Fl_n(\FF_q)| = [n]_q!$ and $|C_w| = q^{\ell(w)}$ for every $w$. \OK{} for
$(n,q) \in \{(2,2),(2,3),(3,2),(3,3),(4,2)\}$, computed from the rank matrix of
Definition~\ref{def:flag} rather than assumed.
\end{proposition}

\begin{center}
\small
\begin{tabular}{rrrrr}
\toprule
$n$ & $q$ & flags & $[n]_q!$ & cells \\
\midrule
$2$ & $2$ & $3$   & $3$   & $2$  \\
$2$ & $3$ & $4$   & $4$   & $2$  \\
$3$ & $2$ & $21$  & $21$  & $6$  \\
$3$ & $3$ & $52$  & $52$  & $6$  \\
$4$ & $2$ & $315$ & $315$ & $24$ \\
\bottomrule
\end{tabular}
\end{center}

\section{The operator, built from the geometry}

\begin{definition}[the fibration, and its operator]\label{def:op}
Let $\pi_i$ forget the $i$-th subspace of a flag. Define $S_i$ on functions on
$\Fl_n(\FF_q)$ by pushforward followed by pullback:
\[
  (S_i f)(F) \;=\; \sum_{\pi_i(F') = \pi_i(F)} f(F').
\]
\end{definition}

Nothing about Schubert cells enters. Two things then have to be checked rather
than assumed: that the fibre is what it should be, and that $S_i$ descends to the
cell module at all.

\begin{proposition}\label{prop:desc}
Every fibre of $\pi_i$ has exactly $q+1 = |\mathbb{P}^1(\FF_q)|$ elements, and
$S_i$ applied to a cell indicator is constant on cells, hence descends to the free
module on $S_n$. \OK{}, $5/5$ cases.
\end{proposition}

\begin{example}[$n = 3$, $q = 2$, the operator read off]\label{ex:rule}
In the basis of cell classes,
\[
  S_1[w] \;=\;
  \begin{cases}
    [w] + [w s_1], & \ell(w s_1) > \ell(w),\\[0.2em]
    q[w] + q[w s_1], & \ell(w s_1) < \ell(w).
  \end{cases}
\]
For instance $S_1[123] = [123] + [213]$ and $S_1[213] = 2[123] + 2[213]$.
\end{example}

\section{What it satisfies, and what it does not}

\begin{proposition}\label{prop:rels}
With $T_i := S_i - I$:
\begin{enumerate}[nosep,label=(\alph*)]
\item $S_i^2 = (q+1) S_i$, so $S_i/(q+1)$ is idempotent;
\item $T_i^2 = (q-1) T_i + q I$, the Hecke quadratic;
\item $T_i T_j T_i = T_j T_i T_j$ for $|i - j| = 1$, and $T_i T_j = T_j T_i$ for
      $|i - j| \ge 2$;
\item $T_i^2 \ne 0$ in every case tested.
\end{enumerate}
\OK{} for (a)--(c), $5/5$ cases, all $24$ cells at $n = 4$.
\end{proposition}

Part (c) is the half of Conjecture~\ref{conj:nil} that survives, and it survives
exactly as predicted: the braid relations come from relabelling positions, need
no new machinery, and are what makes $\Schub_w$ independent of the reduced word.

Part (d) is the half that fails. It is worth being precise about how badly.

\begin{theorem}[no divided difference at finite $q$]\label{thm:nonil}
Fix $i$ and let $A_i$ be the subalgebra generated by $T_i$. Then
$A_i = \operatorname{span}\{I, T_i\}$, and the only $u \in A_i$ with $u^2 = 0$ is
$u = 0$, unless $q = -1$.
\end{theorem}

\begin{proof}
$A_i = \operatorname{span}\{I, T_i\}$ because Proposition~\ref{prop:rels}(b)
expresses $T_i^2$ in that span, so the span is closed under multiplication;
\texttt{flag.py} confirms $T_i^2$ and $T_i^3$ lie in it for all cases tested.
Write $u = aT_i + bI$. Then
\[
  u^2 = a^2 T_i^2 + 2ab\,T_i + b^2 I
      = \bigl[a^2(q-1) + 2ab\bigr] T_i + \bigl[a^2 q + b^2\bigr] I .
\]
$T_i$ and $I$ are linearly independent, so $u^2 = 0$ forces both brackets to
vanish. If $a = 0$ the second gives $b = 0$. If $a \ne 0$ the first gives
$b = a(1-q)/2$, and substituting into the second,
\[
  a^2 q + \tfrac{1}{4}a^2 (1-q)^2 = 0
  \;\Longleftrightarrow\;
  4q + (1-q)^2 = 0
  \;\Longleftrightarrow\;
  (q+1)^2 = 0 . \qedhere
\]
\end{proof}

\begin{corollary}
No nilpotent divided difference exists on the flag container over any finite
field, since $q \ge 2$.
\end{corollary}

\begin{remark}[the scope of the theorem, stated plainly]
The theorem is about operators \emph{attached to a single index}, which is what
a $\partial_i$ is. It does not claim the full algebra has no nilpotent elements
--- $H_n(q)$ is a matrix algebra for generic $q$ and has many. What it claims is
that none of them is a $\partial_i$, and the two-dimensionality of $A_i$ is what
makes the affine search exhaustive rather than a sample. A box search over
$|a|, |b| \le 6$ independently returns nothing in all five cases.
\end{remark}

\begin{refuted}\label{ref:nil}
Conjecture~\ref{conj:nil} is \NO{} in its nilCoxeter half. The flag container
produces $H_n(q)$, whose $q \to 0$ degeneration is the $0$-Hecke algebra of
\emph{idempotent} operators, not the nilCoxeter algebra of nilpotent ones.
\end{refuted}

\begin{center}
\small
\begin{tabular}{lll}
\toprule
operator & relation & generates \\
\midrule
$S_i/(q+1)$      & idempotent        & Demazure operators, \textbf{key polynomials} \\
$T_i = S_i - I$  & Hecke quadratic   & $H_n(q)$, Hall--Littlewood, Green \\
$\partial_i$     & nilpotent         & \textbf{Schubert polynomials} --- not reachable here \\
\bottomrule
\end{tabular}
\end{center}

\section{What the Schubert basis is, at container level}

The refutation is not a dead end; it relocates the object.

\begin{proposition}
The flag container is $\sum_{w \in S_n} X^{\ell(w)}$, so a Schubert class is a
\emph{shape} of the flag container together with its position space
$\FF_q^{\ell(w)}$, and the one-variable shadow of $\Schub_w$ is $x^{\ell(w)}$.
Decategorifying the whole container gives the Poincar\'e polynomial
$\sum_w x^{\ell(w)}$, which at $x = q$ is $[n]_q!$.
\end{proposition}

This says why the divided differences were never going to appear. The
multivariable $\Schub_w(x_1, \ldots, x_n)$ records the \emph{torus weights} on
positions; a point count records only how many positions there are. The two
decategorifications of the same cell decomposition are
\[
  \text{count } \FF_q\text{-points} \;\longrightarrow\; H_n(q), \ [n]_q!
  \qquad\text{versus}\qquad
  \text{Chow ring} \;\longrightarrow\; \partial_i, \ \Schub_w ,
\]
and \texttt{q-container.pdf} is built on the first.

\begin{question}[the row the Targets table is missing]\label{q:graded}
Grade the positions of the flag container by the $T$-action, so that a position
carries a weight rather than only contributing to a dimension. Does the
multigraded flag container have $\Schub_w(x_1,\ldots,x_n)$ as its shape
polynomial, and do the divided differences appear as operations on the grading?
\end{question}

\begin{remark}[$q = -1$]
Theorem~\ref{thm:nonil} fails exactly at $q = -1$, which is not a prime power and
so lies outside the theory rather than inside it at a special point. We record
the value and claim nothing from it.
\end{remark}

\section{Status of every claim}

\begin{center}
\small
\begin{tabular}{clll}
\toprule
\# & claim & prior & verdict \\
\midrule
$5$ & Bruhat decomposition, by enumeration & --- & \OK{}, $5/5$ \\
$6$ & braid relations hold at container level & $70\%$ & \OK{}, $n \le 4$ \\
$7$ & $\partial_i^2 = 0$ at container level & $60\%$ & \NO{}, Thm~\ref{thm:nonil} \\
$8$ & multigraded flag container gives $\Schub_w(x)$ & new & open, Q~\ref{q:graded} \\
\bottomrule
\end{tabular}
\end{center}

Independently of the container side, all $6$ and all $24$ Schubert polynomials
for $S_3$ and $S_4$ were computed in $\ZZ[x]$ by divided differences from
$x^\delta$, with agreement across reduced words asserted at every step; this is
the control, and it confirms $\Schub_{132} = x_1 + x_2$ and
$\Schub_{321} = x_1^2 x_2$ against the standard tables.

\medskip
\noindent\textbf{Reproducing.} \texttt{python3 code/flag.py}. Exact enumeration
over $\FF_2$ and $\FF_3$ for $n \le 4$; no sampling anywhere.

\newcommand{\secindresref}{}
\end{document}
